% \begin{abstract}
% % 需要修改 要求4-6句
% % To enhance the reasoning capabilities of large language models (LLMs), self-consistency has gained significant popularity by combining multiple sampling with majority voting. However, the state-of-the-art self-consistency approaches consume substantial computational resources and lead to significant additional time costs due to the multiple sampling. This prevents its full potential from being realized in scenarios where computational resources are critical.
% % To improve the inference efficiency, this paper introduces \textit{path-consistency}, a method that leverages the confidence of answers generated in earlier branches to identify the prefix of the most promising path. By dynamically guiding the generation of subsequent branches based on this prefix, the \textit{path-consistency} mitigates both the errors and redundancies from random or less useful sampling in self-consistency. As a result, it can significantly accelerate the inference process by reducing the number of tokens generated. Our extensive empirical evaluation shows that the \textit{path-consistency} achieves significant acceleration in inference latency ranging from $7.8\%$ to $40.5\%$, while maintaining or even improving task accuracy across different datasets, including mathematical reasoning, common sense reasoning, symbolic reasoning, and code generation.

% To enhance the reasoning capabilities of large language models (LLMs), self-consistency has become a popular approach, combining multiple samplings with majority voting. However, current methods are computationally expensive and time-consuming due to the need for numerous samplings. To address this, this paper introduces \textit{path-consistency}, which leverages the confidence of earlier-generated answers to identify the most promising prefix and guide the generation of subsequent branches. By dynamically guiding the generation of subsequent branches based on this prefix, \textit{path-consistency} mitigates both the errors and redundancies from random or less useful sampling in self-consistency. This approach reduces errors and redundancies from random sampling, significantly accelerating inference by minimizing token consumption. Our extensive empirical results demonstrate that \textit{path-consistency} improves inference latency by up to 40.5\%, while maintaining task accuracy across various tasks, including mathematical reasoning, commonsense reasoning, and symbolic reasoning.
% \end{abstract}

\begin{abstract}
    Large language models increasingly rely on test-time reasoning strategies such as chain-of-thought prompting and self-consistency to improve performance on complex reasoning tasks. However, self-consistency repeatedly samples full reasoning paths, leading to substantial token redundancy and inference latency, especially when many sampled paths are either incorrect or share similar intermediate reasoning. To address this, we propose path-consistency, a prefix-guided self-consistency method that extracts high-confidence partial reasoning paths from early samples and reuses them as prefixes to guide subsequent generations. By reducing the need to generate each branch from scratch, path-consistency improves token efficiency while preserving the diversity benefits of multi-path sampling. Experiments on arithmetic, commonsense, and symbolic reasoning benchmarks show that path-consistency reduces inference latency by up to 40.5\% while largely preserving the accuracy of standard self-consistency. These results suggest that more efficient use of test-time compute can provide a practical alternative to naively increasing the number of sampled reasoning paths.
\end{abstract}

% Recent advances in large language models have increasingly relied on test-time reasoning strategies, such as chain-of-thought prompting and self-consistency, to improve performance on complex reasoning tasks. However, self-consistency scales test-time compute by repeatedly sampling full reasoning paths, which introduces substantial redundancy in token generation and inference latency, especially when many sampled paths are either incorrect or share similar intermediate reasoning. To address this inefficiency, we propose prefix-guided self-consistency, a lightweight inference-time framework that reuses high-confidence partial reasoning paths from earlier samples as prefixes to guide subsequent generations. Instead of independently generating every reasoning path from scratch, our method dynamically extracts promising prefixes based on answer confidence and conditions later branches on these prefixes, thereby reducing redundant token generation while preserving the diversity benefits of self-consistency. The proposed framework is model-agnostic, requires no additional training or model modification, and can be combined with existing sampling-based reasoning methods. Experiments across arithmetic, commonsense, and symbolic reasoning benchmarks show that prefix-guided self-consistency reduces inference latency by up to 40.5% while largely preserving the accuracy benefits of standard self-consistency. These results suggest that more efficient use of test-time compute can provide a practical alternative to naively increasing the number of sampled reasoning paths.